\chapter{Chest Infection}
\index{Chest Infection}

\section*{Definition and Overview}
A chest infection is an infection of the lungs or airways, most commonly caused by viruses or bacteria. The two main types are acute bronchitis, in which the larger airways (bronchi) become inflamed, and pneumonia, in which the infection reaches the air sacs of the lung itself. Chest infections cause cough, often with phlegm, together with fever, breathlessness, and chest discomfort. Many are mild and clear within one to three weeks, but pneumonia can be serious, particularly in older people, young children, and those with long-term illness. Most chest infections are viral and do not require antibiotics; bacterial infections do. Distinguishing the two is an important part of clinical care.

\section*{Epidemiology}
Chest infections are extremely common and account for a large proportion of visits to primary care and hospitals. Acute bronchitis is most often seen in autumn and winter, when respiratory viruses such as influenza, respiratory syncytial virus, and rhinoviruses circulate widely. Pneumonia is a leading infectious cause of death worldwide, especially among older adults, young children, and people with chronic heart, lung, or immune disease. Risk is increased by smoking, alcohol, overcrowding, and conditions that impair coughing or the immune system. Vaccination against influenza, pneumococcus, and other respiratory pathogens has substantially reduced the burden of serious chest infection in many populations.

\section*{Aetiology and Pathophysiology}
\begin{itemize}
    \item \textbf{Viruses:} the most common cause of chest infection; influenza, respiratory syncytial virus, rhinoviruses, adenoviruses, and coronaviruses are frequent agents
    \item \textbf{Bacteria:} \textit{Streptococcus pneumoniae}, \textit{Haemophilus influenzae}, \textit{Mycoplasma pneumoniae}, and \textit{Legionella} are important causes of pneumonia
    \item \textbf{Bronchitis mechanism:} viral infection inflames the lining of the larger airways, causing swelling, excess mucus, and coughing
    \item \textbf{Pneumonia mechanism:} the infection reaches the alveoli, where inflammatory fluid and white cells fill the air spaces, reducing oxygen exchange
    \item \textbf{Transmission:} spread by droplets and aerosols from coughing and sneezing, or by touching contaminated surfaces
    \item \textbf{Host factors:} smoking, chronic lung disease, heart failure, diabetes, and impaired immunity all increase susceptibility and severity
    \item \textbf{Complications pathway:} severe inflammation can progress to respiratory failure, sepsis, and multi-organ dysfunction
\end{itemize}

\section*{Clinical Features}
\begin{itemize}
    \item \textbf{Cough:} the dominant symptom, often productive of green, yellow, or rusty-coloured phlegm
    \item \textbf{Fever:} raised temperature, sometimes with chills and sweats
    \item \textbf{Breathlessness:} difficulty breathing or rapid breathing, particularly with pneumonia
    \item \textbf{Chest pain:} often pleuritic -- sharp and worse on breathing or coughing
    \item \textbf{General symptoms:} fatigue, muscle aches, headache, and poor appetite
    \item \textbf{Wheeze:} more prominent with bronchitis, especially in people with asthma
    \item \textbf{Severe features:} confusion, bluish lips, low oxygen levels, coughing up blood, and an inability to keep fluids down
    \item In older adults and infants, the main signs may be confusion, lethargy, poor feeding, or rapid breathing rather than a prominent cough
\end{itemize}

\section*{Differential Diagnosis}
\begin{itemize}
    \item Asthma, COPD, and other causes of airway obstruction and wheeze
    \item Heart failure with fluid in the lungs
    \item Pulmonary embolism, which can cause breathlessness and chest pain
    \item Bronchiectasis, where chronic cough and phlegm are longstanding
    \item Tuberculosis, especially with prolonged cough, weight loss, and night sweats
    \item Lung cancer, which can present with cough, haemoptysis, and recurrent infection
    \item Gastro-oesophageal reflux or post-nasal drip causing chronic cough
    \item COVID-19 and other viral respiratory infections
\end{itemize}

\section*{When to Seek Medical Care}
See a doctor if a cough lasts longer than three weeks, if you are coughing up blood, if you have a high fever that does not settle, or if you have underlying lung, heart, or immune disease and develop symptoms of chest infection.

\medskip
\textbf{Contact your local emergency services for:}
\begin{itemize}
    \item Severe difficulty breathing, gasping, or breathing very fast
    \item Bluish or grey lips, face, or fingernails
    \item Chest pain, confusion, or drowsiness
    \item Coughing up blood
    \item Signs of sepsis, such as a mottled rash, cold extremities, or reduced urine output
    \item A baby who is unwell, not feeding, or breathing rapidly
\end{itemize}

\section*{Diagnosis and Investigations}
\begin{itemize}
    \item \textbf{History:} duration and nature of cough, fever, breathlessness, smoking history, and underlying health conditions
    \item \textbf{Examination:} measuring temperature, oxygen saturation, respiratory rate, and listening to the chest for crackles or wheeze
    \item \textbf{Pulse oximetry:} a simple test of oxygen levels in the blood, used to judge severity
    \item \textbf{Chest X-ray:} distinguishes pneumonia from bronchitis and may reveal complications such as fluid or lung collapse
    \item \textbf{Blood tests:} inflammatory markers and full blood count help indicate a bacterial infection; blood cultures may be taken in hospital
    \item \textbf{Sputum culture:} may guide antibiotic choice in severe or non-improving cases
    \item \textbf{Viral testing:} nasal or throat swabs (for example, for influenza or COVID-19) inform management in hospital
    \item \textbf{Further imaging or bronchoscopy:} considered when infection is slow to resolve or another diagnosis is suspected
\end{itemize}

\section*{Management}
\begin{itemize}
    \item \textbf{Self-care for mild infection:} rest, plenty of fluids, and simple pain and fever relief such as paracetamol
    \item \textbf{Antibiotics:} prescribed only when a bacterial cause is likely or confirmed; most chest infections are viral and antibiotics will not help
    \item \textbf{Inhaled bronchodilators:} reliever inhalers can ease wheeze and coughing in those with underlying airways disease
    \item \textbf{Hospital care:} oxygen, intravenous fluids, and intravenous antibiotics for severe or complicated pneumonia
    \item \textbf{Supportive measures:} physiotherapy and breathing exercises to help clear secretions
    \item \textbf{Follow-up:} a chest X-ray may be repeated after several weeks to confirm pneumonia has fully cleared, especially in smokers
    \item \textbf{Antivirals:} specific treatments may be offered for influenza in high-risk people within a suitable timeframe
\end{itemize}

\section*{Complications}
\begin{itemize}
    \item Pleural effusion -- fluid collecting around the lung
    \item Empyema -- infected pus in the space around the lung, requiring drainage
    \item Lung abscess, in severe necrotising infections
    \item Respiratory failure requiring oxygen or ventilatory support
    \item Sepsis and septic shock
    \item Delirium or worsening of underlying conditions, especially in older adults
    \item Prolonged cough and fatigue lasting several weeks
    \item Spread of infection to other sites, including the bloodstream and meninges
\end{itemize}

\section*{Prognosis and Outcomes}
Most people with a chest infection recover fully within one to three weeks, although the cough can linger for several weeks afterwards as the airways heal. Viral bronchitis usually resolves without specific treatment. Pneumonia is more serious: most cases in healthy adults respond well to treatment, but the risk of severe illness and death rises with age and underlying disease. Prompt medical assessment and early treatment substantially improve outcomes, and most complications are treatable. People who smoke, and those with chronic lung or heart disease, should be followed up after pneumonia to confirm complete recovery.

\section*{Prevention and Self-Care}
\begin{itemize}
    \item Keep vaccinations up to date, including influenza, pneumococcal, and COVID-19 vaccines
    \item Stop smoking and avoid smoky environments
    \item Wash hands regularly and cover coughs and sneezes
    \item Stay home and rest when unwell with a respiratory infection
    \item Drink plenty of fluids and take simple pain relief when symptoms develop
    \item Avoid antibiotics unless prescribed, to reduce resistance
    \item Manage underlying conditions such as asthma, diabetes, and heart failure well
    \item Protect vulnerable people by avoiding contact when you have a cold or flu-like illness
\end{itemize}

\section*{Sources and Further Reading}
The following reputable sources provide further information on chest infections:
\begin{itemize}
    \item NHS. \textit{Health A--Z: chest infection information.} https://www.nhs.uk/conditions/
    \item Mayo Clinic. \textit{Diseases and conditions library.} https://www.mayoclinic.org/
    \item NICE. \textit{Clinical guidelines on pneumonia and respiratory infections.} https://www.nice.org.uk/
    \item MSD Manual. \textit{Professional edition: pulmonary disorders.} https://www.msdmanuals.com/
    \item Centers for Disease Control and Prevention. \textit{Pneumonia and respiratory infection resources.} https://www.cdc.gov/
    \item World Health Organization. \textit{Pneumonia fact sheets.} https://www.who.int/
\end{itemize}
